
\vspace{5mm}
\begin{ledgroup}
\footnotesize
\pstart
\noindent
Die vorliegenden, auf zwei Bogen überlieferten drei Stücke gehen auf Leibnizens Lektüre von 
\indexN{\textls{Saccheri}, Girolamo 1667\textendash1733}Girolamo Saccheris 
\cite{00541}\textit{Neo-statica} (Mailand 1708) im Winter 1712/1713 
zurück.
\pend
\pstart
Bei N.\;\ref{60123} handelt es sich um Auszüge mit Bemerkungen aus 
\indexN{\textls{Saccheri}, Girolamo 1667\textendash1733}Saccheris \cite{00541}Werk, 
die die ersten drei Seiten des Bogens Bl.~3\textendash4 einnehmen.
Das auf der letzten Seite desselben Bogens, und aller Wahrscheinlichkeit nach im Anschluss an N.\;\ref{60123}, verfasste Konzept N.\;\ref{60124} 
greift die Themen der zusammengesetzten Bewegung und der Proportion parazentrischer Impulse besonders hervor, wobei Leibniz
\indexN{\textls{Saccheri}, Girolamo 1667\textendash1733}Saccheris Überlegungen 
eigene Theorien gegenüberstellt.
N.\;\ref{60124} diente wiederum als Vorlage für das ausführlichere Konzept N.\;\ref{60122}; 
der (von Leibniz nachträglich vollständig gestrichene) Text von N.\;\ref{60124} 
wurde weitgehend und mit teils beträchtlichen wörtlichen Übereinstimmungen
in N.\;\ref{60122} übernommen
(vgl. bspw. die Passagen \refpassage{35_10_16_004_4a}{35_10_16_004_4b} und \refpassageohne{35_10_16_004_5b}{35_10_16_004_5b} von N.\;\ref{60124} mit \refpassage{35_10_16_001-002_7a}{35_10_16_001-002_7b} und \refpassageohne{35_10_16_001-002_8a}{35_10_16_001-002_8b} in N.\;\ref{60122}).
Eine genauere Datierung ermöglichen einige von Leibniz in N.\;\ref{60122} gemachte Angaben.
\pend
\pstart
Einzig das zuletzt entstandene Konzept N.\;\ref{60122} hat Leibniz datiert: \glqq Viennae\indexO{Wien} mense Februario 1713\grqq\ (\refpassage{35_10_16_001-002_6a}{35_10_16_001-002_6a}); 
dieses Datum bildet einen Terminus ante quem für den gesamten Textkomplex.
Zudem bemerkt er an derselben Stelle, dass  
\indexN{\textls{Saccheri}, Girolamo 1667\textendash1733}Saccheris 
\cite{00541}\textit{Neo-statica} zwar um 1708 erschienen sei,
er sie aber \glqq erst jetzt\grqq\ habe lesen können (\glqq jam Mediolani\indexO{Mailand} A.\,D.~1708 edita, sed nunc demum a R\textsuperscript{mo} Patre Granello\indexN{\textls{Granelli} (Granellus), Carlo 1671\textendash1739} 
mihi communicata\grqq, \refpassage{35_10_16_001-002_6a}{35_10_16_001-002_6b}).
Der Zeitpunkt der (allen drei Stücken zugrundeliegenden, in den Auszügen N.\;\ref{60123} festgehaltenen) 
Lektüre der \cite{00541}\textit{Neo-statica} lässt sich genauer bestimmen,
da Leibniz wahrscheinlich Mitte Dezember 1712 in Wien ankam (\cite{01236}\textit{Chronik}, S.\;233).
Im Zeitraum Mitte Dezember 1712 bis Februar 1713 (der mit der Angabe "nunc demum" gut übereinstimmt) müsste Leibniz von einem "Patre Granello" das Werk erhalten haben.
Bei Letztgenanntem handelt es sich 
um den Jesuiten \indexN{\textls{Granelli} (Granellus), Carlo 1671\textendash1739}Carlo Granelli, 
Historiker und seit 1710 Beichtvater der früheren
\indexN{\textls{Braunschweig-Lüneburg}, Wilhelmine Amalie von, Kaiserin des HRR 1705\textendash1711 (1673\textendash1742)}Kaiserin Wilhelmine Amalie 
in Wien\indexO{Wien}
(vgl. \textsc{A.~Cito}, \textit{Vita, e virtù dell'imperadrice Guglielmina Amalia}, Wien~1744, S.~\lbrack XVI\rbrack). 
Granellis Exemplar der \textit{Neo-statica} ist in der Universitätsbibliothek Wien aufbewahrt (I-249050; zeitgenössisches Provenienzvermerk: \glqq Ex Dono R.\;P.\textsuperscript{is} Caroli Granelli S.\,J. 1732\grqq; das Buch enthält keine Anstreichungen oder Marginalien.) 
Aus diesen Umständen ergibt sich also die vorgeschlagene Datierung: Mitte Dezember 1712 bis Februar 1713.
\pend
\pstart
Ein Hinweis in Leibnizens
\cite{02124}Brief an \indexN{\textls{Hermann}, Jacob 1678\textendash1733}J.~Hermann 
vom 1.~Februar 1713
(LBr~396, Bl.~98\textendash99; ediert in \cite{00150}\textit{LOD}~III, S.\;536\textendash538)
bestätigt seine Kenntnis der \cite{00541}\textit{Neo-statica} zu jenem Zeitpunkt.
Leibniz schreibt: \glqq Vidistine Sacchieri jesuitae \indexO{Pavia}Ticinensis \textit{Neostaticam}, ubi supponitur corpora gravia habere directiones non parallelas, sed concurrentes in centro. Valde virum laudant, et magni ingenii esse ajunt\grqq\ (Bl.~99~v\textsuperscript{o}).
Noch einige Monate später erwähnt Leibniz in seinem
\cite{02123}Brief an \indexN{\textls{Des Bosses}, Barthélemy 1638\textendash1738}Des~Bosses 
vom 23. August 1713 (vgl. die Abfertigung, Paris \textit{BnF} Latin~10355, Bl.~6\textendash7; ediert in \cite{02039}\textit{LPG}~II, S.\;481\textendash483)
\indexN{\textls{Saccheri}, Girolamo 1667\textendash1733}Saccheris Werk:
\glqq Vidi quaedam ingeniosa admodum scripta a R.\,P. Saccherio Mathematico Ticinensi, ex Ordine vestro, in lucem edita, etsi sententias quasdam mathematicas foveat, quas non omnino probare possum\grqq\ (Bl.~7~v\textsuperscript{o}; bei \cite{02039}Gerhardt, S.\;482, inkorrekt als \glqq R.\,P.~Sacorierio\grqq\ transkribiert).
\pend
\pstart
Leibniz hat einige seiner Kommentare in N.\;\ref{60124}
zunächst in der zweiten Person formuliert (\glqq ponamus tecum...\grqq\ u.ä.),
womöglich ein Indiz seiner Absicht, sie \indexN{\textls{Saccheri}, Girolamo 1667\textendash1733}Saccheri mitzuteilen;
doch hat er diese Ausdrucksweise gleich verworfen
(siehe die Streichungen zu \refnote{FnB:LH_35_10_16_004_esto-tecum}, \refnote{FnB:LH_35_10_16_004_consentire-tuae} und \refnote{FnB:LH_35_10_16_004_ponamus-tecum}).
Ein Briefwechsel zwischen Leibniz und \indexN{\textls{Saccheri}, Girolamo 1667\textendash1733}Saccheri ist nicht bekannt.
\pend
\end{ledgroup}
